\section{问题建模}

\label{sec:problem-formulation}

本文研究强约束三维低空无人机路径规划问题。给定三维飞行空间、起点、终点、障碍物、禁飞区和高度约束，目标是在安全可行前提下生成平滑路径。与仅最小化长度的规划不同，低空飞行还需兼顾能耗、风险暴露和转弯平滑性，并同时满足高度与空间边界约束。因此，本文将路径规划写成带约束的连续优化问题，并采用可行性优先准则比较候选解。

\subsection{三维路径表示}

设无人机起点和终点分别为 $\mathbf{s}=[x_s,y_s,z_s]^\mathsf{T}$ 和 $\mathbf{g}=[x_g,y_g,z_g]^\mathsf{T}$。优化变量为 $K$ 个内部控制点：
\begin{equation}
    \mathbf{X}=
    [\mathbf{q}_1^\mathsf{T},\mathbf{q}_2^\mathsf{T},\ldots,\mathbf{q}_K^\mathsf{T}]^\mathsf{T}\in\mathbb{R}^{3K}
    \quad
    \mathbf{q}_k=[x_k,y_k,z_k]^\mathsf{T}
    \label{eq:path-control-vector}
\end{equation}
完整控制点序列为
\begin{equation}
    \mathcal{Q}=\{\mathbf{s},\mathbf{q}_1,\ldots,\mathbf{q}_K,\mathbf{g}\}
    \label{eq:control-sequence}
\end{equation}
本文使用 open-uniform 三次 B 样条由控制点序列构造连续路径 $\mathbf{p}(u)$，$u\in[0,1]$。实际评价中取 $M$ 个均匀采样点：
\begin{equation}
    \begin{aligned}
    \mathcal{P}(\mathbf{X})&=\{\mathbf{p}_m\}_{m=1}^{M}\\
    \mathbf{p}_m&=\mathbf{p}\!\left(\frac{m-1}{M-1}\right)=[x_m,y_m,z_m]^\mathsf{T}
    \end{aligned}
    \label{eq:path-sampling}
\end{equation}
因此 $\mathbf p_1=\mathbf s$ 且 $\mathbf p_M=\mathbf g$。
路径长度、风险、碰撞、禁飞区、高度和转弯约束均基于采样点序列计算，从而在较低控制点维度下统一处理多种空间约束。

\subsection{优化目标}

对任一路径 $\mathcal{P}(\mathbf{X})$，本文采用综合路径质量代理目标
\begin{equation}
    J(\mathbf{X}) =
    w_L L + w_E E + w_R R + w_S S + P_c
    \label{eq:objective}
\end{equation}
其中 $L,E,R,S$ 分别表示长度、能耗、风险和平滑性，$w_L,w_E,w_R,w_S$ 为对应权重。惩罚项 $P_c$ 写为
\begin{equation}
    P_c = \lambda_{\mathrm{obs}} C_{\mathrm{obs}}
    + \lambda_{\mathrm{nfz}} C_{\mathrm{nfz}}
    + \lambda_{\mathrm{curv}} C_{\mathrm{curv}}
    + \lambda_{\mathrm{alt}} C_{\mathrm{alt}}
    \label{eq:penalty}
\end{equation}
式 \eqref{eq:objective} 中各目标项基于采样序列定义。令相邻采样点之间的段向量和段长分别为
\begin{equation}
    \Delta \mathbf{p}_m=\mathbf{p}_{m+1}-\mathbf{p}_m,\quad
    \Delta l_m=\|\Delta \mathbf{p}_m\|_2+\epsilon_l
    \label{eq:segment-definition}
\end{equation}
其中 $m=1,\ldots,M-1$，$\epsilon_l$ 为防止零长度除法的极小常数。路径长度项定义为
\begin{equation}
    L=\sum_{m=1}^{M-1}\Delta l_m
    \label{eq:length-term}
\end{equation}
设 $\mathbf{w}_m$ 为第 $m$ 段起点处的风场向量，$\widehat{\mathbf d}_m=\Delta\mathbf p_m/\Delta l_m$ 为带除零保护的段方向，$\widehat{\mathbf w}_m=\mathbf w_m/(\|\mathbf w_m\|_2+\epsilon_l)$ 为带除零保护的风向方向。风阻相关项写为
\begin{equation}
    \Phi_m=\|\mathbf w_m\|_2
    \left(1-\widehat{\mathbf d}_m^\mathsf{T}\widehat{\mathbf w}_m\right)
    \label{eq:wind-alignment}
\end{equation}
能耗项由水平/空间飞行距离、垂向变化和逆风暴露共同构成：
\begin{equation}
    E=\sum_{m=1}^{M-1}
    \left(
    a_1\Delta l_m+a_2|\Delta z_m|+a_3\Phi_m
    \right)
    \label{eq:energy-term}
\end{equation}
其中 $\Delta z_m=z_{m+1}-z_m$，$a_1,a_2,a_3$ 为能耗系数。若 $r_{\mathrm{obj}}(\mathbf p_m)$ 为评价器使用的采样点风险密度，则风险项定义为
\begin{equation}
    R=\sum_{m=1}^{M-1}r_{\mathrm{obj}}(\mathbf p_m)\Delta l_m
    \label{eq:risk-term}
\end{equation}
平滑性项定义为
\begin{equation}
    S=\sum_{m=2}^{M-1}\|\mathbf p_{m+1}-2\mathbf p_m+\mathbf p_{m-1}\|_2^2
    \label{eq:smoothness-term}
\end{equation}
式 \eqref{eq:energy-term} 所定义的项为固定采样设置下的路径质量代理，不解释为严格飞行动力学或真实能耗模型；硬约束违反由下文单独判断。

\subsection{约束条件}

\paragraph{空间边界约束}
设 $\mathbf l_0,\mathbf u_0\in\mathbb R^3$ 为单个控制点的下界和上界，并令 $\mathbf l=[\mathbf l_0^\mathsf T,\ldots,\mathbf l_0^\mathsf T]^\mathsf T$、$\mathbf u=[\mathbf u_0^\mathsf T,\ldots,\mathbf u_0^\mathsf T]^\mathsf T$。决策空间为 $\Omega=\{\mathbf X\in\mathbb R^{3K}:\mathbf l\leq\mathbf X\leq\mathbf u\}$；每个内部控制点和采样点均应位于对应的三维规划空间范围内，即 $\mathbf l_0\leq \mathbf q_k\leq\mathbf u_0$，$k=1,\ldots,K$。
候选控制点越界时先通过 $\Pi_{\Omega}(\cdot)$ 投影回规划盒内。由于本文采用的 open-uniform B 样条基函数非负且和为 $1$，路径点可写成控制点的凸组合；若起点、终点和内部控制点均位于同一规划盒内，则采样路径不会越出该盒。参考路径中的边界偏好仅提供软引导，不构成目标函数项或额外硬边界判定。

\paragraph{障碍物约束}
设障碍物集合为 $\mathcal{O}$。采样点不得进入障碍物，即 $\mathbf{p}_m \notin \mathcal{O}$，$m=1,\ldots,M$。
障碍物违反量按进入障碍物的采样点数量累计：
\begin{equation}
    C_{\mathrm{obs}}=
    \sum_{m=1}^{M}\mathbb{I}(\mathbf p_m\in\mathcal O)
    \label{eq:obs-violation}
\end{equation}
其中 $\mathbb{I}(\cdot)$ 为指示函数。CVF 可在候选生成中使用近障碍缓冲压力，但不改变式 \eqref{eq:obs-violation} 的可行性判定。

\paragraph{禁飞区约束}
设禁飞区集合为 $\mathcal{Z}$，其中 $\mathcal{Z}$ 表示已合并安全缓冲后的禁飞区集合。路径采样点不得进入禁飞柱体或其安全缓冲范围，即 $\mathbf{p}_m \notin \mathcal{Z}$，$m=1,\ldots,M$。
禁飞区违反量同样按采样点进入禁飞区的次数统计：
\begin{equation}
    C_{\mathrm{nfz}}=
    \sum_{m=1}^{M}\mathbb{I}(\mathbf p_m\in\mathcal Z)
    \label{eq:nfz-violation}
\end{equation}

\paragraph{高度约束}
无人机飞行高度需满足场景给定的下界和上界，即 $z_{\min} \leq z_m \leq z_{\max}$，$m=1,\ldots,M$。
若路径低于 $z_{\min}$ 或高于 $z_{\max}$，对应超限距离累计为
\begin{equation}
    C_{\mathrm{alt}}=\sum_{m=1}^{M}
    \left[\max(0,z_{\min}-z_m)+\max(0,z_m-z_{\max})\right]
    \label{eq:altitude-violation}
\end{equation}
不同场景可采用不同高度上界和参考高度 $z_{\mathrm{ref}}$，以反映低空走廊收紧要求；$z_{\mathrm{ref}}$ 仅用于初始化和 CVF 的软引导，不参与适应度函数。

\paragraph{转弯约束}
令 $\mathbf v_m=\mathbf p_{m+1}-\mathbf p_m$。连续路径段形成的转弯角定义为
\begin{equation}
    \begin{aligned}
    \theta_m=
    \arccos
    \left[
    \mathrm{clip}\left(
    \frac{\mathbf v_{m-1}^\mathsf{T}\mathbf v_m}
    {\|\mathbf v_{m-1}\|_2\|\mathbf v_m\|_2+\epsilon_l},
    -1,1
    \right)
    \right]\\
    m=2,\ldots,M-1
    \end{aligned}
    \label{eq:turning-angle}
\end{equation}
其中 $\mathrm{clip}(x,-1,1)$ 表示将标量 $x$ 截断到 $[-1,1]$，用于避免数值舍入导致反余弦输入越界。
路径转弯角不应超过允许上限 $\theta_{\max}$，即 $\theta_m \leq \theta_{\max}$，$m=2,\ldots,M-1$；超出部分累计为
\begin{equation}
    C_{\mathrm{curv}}=\sum_{m=2}^{M-1}\max(0,\theta_m-\theta_{\max}),
    \label{eq:curvature-violation}
\end{equation}
从而将急转弯或局部折线化路径识别为不可行候选。

\subsection{可行性准则}

为避免不同量纲约束被误解为同一物理量，本文首先将约束违反写成非负分量向量
\begin{equation}
    \mathbf c(\mathbf X)=
    \left[
    C_{\mathrm{obs}},
    C_{\mathrm{nfz}},
    C_{\mathrm{curv}},
    C_{\mathrm{alt}}
    \right]^\mathsf T
    \label{eq:constraint-vector}
\end{equation}
本文以四个非负分量之和定义总违反分数
\begin{equation}
    V(\mathbf{X}) =
    C_{\mathrm{obs}} + C_{\mathrm{nfz}} + C_{\mathrm{curv}} + C_{\mathrm{alt}}
    \label{eq:total-violation}
\end{equation}
采用总违反分数阈值判定可行性，即当且仅当 $V(\mathbf{X})\leq\varepsilon_v$ 时路径可行，其中 $\varepsilon_v=10^{-10}$ 仅用于抑制浮点误差。各分量非负，因而该判定仅表示固定采样点和离散转弯角满足数值容差，而非采样点间连续轨迹的碰撞认证。$V$ 是依赖地图单位、采样数和弧度制的经验聚合分数，仅用于不可行候选的排序和阈值判断。

候选解比较采用 Deb 可行性优先规则 \citep{Deb2000ConstraintHandling}。对两个候选解 $\mathbf{X}_a$ 和 $\mathbf{X}_b$，若仅一个解可行，则可行解优先；若二者均可行，则含惩罚综合适应度 $J$ 更小者优先；若二者均不可行，则 $V$ 更小者优先；若 $V$ 相同，再以 $J$ 作为次级判据。该规则使不可行候选的排序倾向于降低总违反分数，并在可行候选之间比较综合适应度。第 \ref{sec:method} 节将进一步说明 CVF-AE 如何利用这一可行性信息构造约束可行性方向。
